\section{Data Engineering, Validity, and Network Construction}
\label{sec:data_methodology}

\subsection{Empirical Setting, Data Ingestion, and Boundary Specification}
To systematically analyze the formation, persistence, and dissolution of Eurovision voting blocs over a half-century, we construct a longitudinal panel of country-to-country voting networks spanning 1975 to 2024. Primary contest results and dyadic voting breakdowns from 1975 to 2023 are extracted from the Spijkervet Eurovision Dataset \citep{spijkervet2023}, while 2024 results are harvested from the curated \texttt{josago97/EurovisionDataset} corpus \citep{josago2025} and cross-validated against official European Broadcasting Union (EBU) archives. Dyadic geographic contiguity and official linguistic alignments are retrieved from the CEPII GeoDist database \citep{mayer2011}.

Following network analysis foundational principles, our boundary specification adheres strictly to a \textbf{nominalist} (\textit{etic}) strategy \citep{laumann1983boundary}. Rather than relying on subjective cultural affiliations (\textit{emic} self-recognition), network boundaries are objectively defined by the formal institutional criteria established by the EBU: any country broadcasting the Eurovision Grand Final whose national jury or tele-voting electorate was accredited to award official points in that year's final. To ensure strict longitudinal comparability and eliminate structural distortion, our analysis is restricted exclusively to the \textbf{Grand Final} rounds. This design choice carries major implications for construct validity.

\subsection{Construct Validity: Mitigating Omission and Commission Errors}
In sociocentric network research, evaluating construct validity requires assessing whether observed ties and nodes faithfully capture the underlying relational phenomena without introducing systematic bias:
\begin{itemize}
    \item \textbf{Omission Errors}: A critical structural vulnerability in Eurovision data stems from the introduction of semi-finals in 2004. In the semi-final rounds, only a subset of nations competes and votes within assigned qualifying heats, creating fragmented subgraphs with artificial structural zeros. However, under EBU governance rules since 2004, all participating nations---including non-qualifying semi-finalists---retain full voting rights during the Saturday Grand Final. Restricting the boundary to Grand Final decisions avoids omitting voting populations, ensuring every sovereign participant maintains an equal structural opportunity to cast preference ties.
    \item \textbf{Commission Errors}: Starting in 2023, the EBU introduced an aggregate ``Rest of the World'' (ROW) public online vote acting as an additional synthetic country. Because ROW does not correspond to a sovereign geographical territory, lacks bilateral dyadic attributes (borders, official languages), and cannot receive votes, its inclusion would represent a severe commission error, distorting in-degree/out-degree balance and community partitions. Consequently, ROW votes were systematically filtered from the network.
\end{itemize}

\subsection{Reliability: Harmonization and Historical Geopolitics}
Ensuring measurement reliability (achieving identical network topographies under replication across disparate archival repositories) required resolving significant formatting discrepancies and geopolitical transitions:
\begin{enumerate}
    \item \textbf{Geopolitical Successions}: The contest's longitudinal scope traverses major geopolitical dissolutions. Historical Yugoslavia (\texttt{YUG}) is tracked across Period 1 (1975--1996), the Federal Republic of Yugoslavia / State Union of Serbia and Montenegro (\texttt{SCG}) across Period 2 (1997--2008), and independent Serbia (\texttt{SRB}) and Montenegro (\texttt{MNE}) across Periods 3 and 4. All ISO country codes were standardized to the ISO 3166-1 alpha-3 standard.
    \item \textbf{Microstates and GeoDist Augmentation}: CEPII GeoDist records historical nation-states under legacy abbreviations (e.g., \texttt{ROM} for Romania, standardized to \texttt{ROU}) and omits microstates such as Monaco (\texttt{MCO}) and San Marino (\texttt{SMR}), as well as transitional entities (\texttt{SCG}). To guarantee complete zero-loss matrix alignment for Quadratic Assignment Procedure (QAP) hypothesis testing, we explicitly augmented CEPII matrices with sovereign border verifications (Monaco contiguity with France; San Marino enclave contiguity with Italy) and official language codings.
\end{enumerate}

\subsection{Institutional Eras and Longitudinal Top-$k$ Network Formulation}
Between 1975 and 2024, the contest underwent four fundamental institutional shifts in governance, voting technology, and scoring mechanisms:
\begin{itemize}
    \item \textbf{Period 1 (1975--1996, Jury Era)}: 22 finals ($N=35$). Regulated solely by closed national expert juries awarding points under the classic 12-10-8--1 scale.
    \item \textbf{Period 2 (1997--2008, Televote Rollout \& Expansion)}: 12 finals ($N=48$). Rapid phase-in of public televoting via telephone and SMS; post-Cold War eastern enlargement; semi-finals introduced in 2004.
    \item \textbf{Period 3 (2009--2015, Combined 50/50 System)}: 7 finals ($N=47$). Re-introduction of juries combined equally (50\% televote + 50\% jury) into a single aggregate rank to mitigate alleged diaspora bias.
    \item \textbf{Period 4 (2016--2024, Separate Points Presentation)}: 8 finals ($N=45$). National juries and public televoters independently distribute separate 12-point scales (combined total used here for longitudinal comparability; 2020 contest was cancelled).
\end{itemize}

\begin{table}[htbp]
\centering
\caption{Eurovision Voting Networks Dataset Summary across Eras (1975--2024)}
\label{tab:dataset_summary}
\small
\begin{tabular}{l l c c c c c c c c}
\hline\hline
Period & Voting System & Years & Finals & $N$ & $E_{\text{weighted}}$ & Mean $W$ & $E_{t \ge 0.5}$ & $E_{t \ge 0.33}$ & Contig. (\%) & ComLang (\%) \\
\hline
P1 & Jury only & 1975--1996 & 22 & 35 & 518 & 0.261 & 74 & 151 & 7.7\% & 5.9\% \\
P2 & Televote rollout & 1997--2008 & 12 & 48 & 615 & 0.387 & 194 & 276 & 8.3\% & 3.2\% \\
P3 & 50/50 Jury and Televote & 2009--2015 & 7 & 47 & 536 & 0.410 & 172 & 296 & 7.9\% & 3.1\% \\
P4 & Separate voting (combined) & 2016--2024 & 8 & 45 & 686 & 0.321 & 148 & 255 & 7.8\% & 3.6\% \\
\hline\hline
\end{tabular}
\begin{tablenotes}
\footnotesize
\item \textit{Note:} $N$ represents active Grand Final participants/voters in the era; $E_{\text{weighted}}$ denotes directed ties where $W(A \to B) > 0$ for top-$k=3$; $E_{t \ge 0.5}$ and $E_{t \ge 0.33}$ are binarized edges at the main ($0.50$) and robustness ($0.33$) persistence thresholds; Contig. and ComLang show the percentage of directed dyads sharing a land border or official language from CEPII GeoDist.
\end{tablenotes}
\end{table}


To overcome the radical rule changes and the annual fluctuation of competing finalists, we adhere to the \textbf{Phenomenon $\to$ Network Measure $\to$ Real-world Interpretation ($A \to B \to C$)} paradigm:
\begin{itemize}
    \item \textbf{Phenomenon ($A$)}: Voting blocs do not manifest as random support; rather, they reflect entrenched, recurring favoritism where a donor state consistently allocates its highest accolades to the same recipient, irrespective of individual song quality.
    \item \textbf{Network Measure ($B$)}: Within each annual contest $t$, we map voting allocations into a binary directed edge $A \to B$ if nation $B$ ranks within the top-$k$ preferences of nation $A$ (primary specification $k=3$, capturing the prestigious 8, 10, and 12-point tiers). Over an institutional era $P$, directed edge weights $W_P(A \to B)$ are normalized by bilateral co-opportunity:
    \begin{equation}
        W_P(A \to B) = \frac{\sum_{t \in P} \mathbb{I}(B \in \text{Top-}k_t(A))}{\sum_{t \in P} \mathbb{I}(A \in \text{Voters}_t \land B \in \text{Finalists}_t)}
    \end{equation}
    To facilitate structural triad census and reciprocity metrics, we binarize the directed adjacency matrix at a strict majority threshold $W_P(A \to B) \ge 0.50$ with a repeat-support constraint $\sum_{t \in P} \mathbb{I}(B \in \text{Top-}k_t(A)) \ge 2$ (primary, \texttt{\{period\}\_adj\_binary\_t05.csv}). This critical constraint eliminates spurious single-year artifact ties generated by peripheral nations participating in only one or two contests during a period (e.g., Morocco in 1980 or newly entered Eastern European states in 1994--1996 with $1/1 = 1.0$). We likewise evaluate a sensitivity threshold $W_P(A \to B) \ge 0.33$ ($\text{top-}k \ge 2$; \texttt{robustness/\{period\}\_adj\_binary\_t033.csv}), alongside alternative configurations for $k=2$, $k=4$, unconstrained thresholds (\texttt{robustness/\{period\}\_adj_binary_t05_unconstrained.csv}), and rank-based weights (\texttt{robustness/\{period\}\_adj_rank.csv}).
    \item \textbf{Real-world Interpretation ($C$)}: As summarized in Table~\ref{tab:dataset_summary}, mean tie weight surges from $0.261$ in the Jury Era (P1) to $0.387$ in P2 and $0.410$ in P3. While geographic contiguity remains structurally invariant across eras ($\sim 7.7\%$--$8.3\%$), the nearly five-fold increase in persistent binary ties surviving at $t \ge 0.5$ with repeat support (from 23 in P1 to 112 in P2 and 101 in P4) directly demonstrates how the democratization of voting via telephone technology activated deep-seated geopolitical affinity networks and migrant diasporas that national professional juries had historically suppressed.
\end{itemize}
